% !PRINTED 170
\weberpage{170}

\begin{center}
§. 52.

\gesperrt{Tschirnhausen-Transformation.}
\end{center}

Wir machen von der Theorie der symmetrischen Functionen und der Resultantenbildung noch eine wichtige Anwendung auf die von \gesperrt{Tschirnhausen}\textsuperscript{1)} herrührende Transformation algebraischer Gleichungen.

Es sei
\weq{1}{f(x) = x^n + a_1 x^{n-1} + a_2 x^{n-2} + \cdots + a_{n-1} x + a_n}
irgend eine Function $n^{\text{ten}}$ Grades, und
\weq{2}{y = \frac{\chi(x)}{\psi(x)}}
irgend eine (ganze oder gebrochene) rationale Function von $x$, bei der wir nur die eine Voraussetzung machen, dass $f(x)$ und $\psi(x)$ \gesperrt{keinen gemeinsamen Theiler haben sollen}. Nach dem Schlusssatz von \S.\,6 können wir dann die Functionen $F(x)$ und $\varphi(x)$, letztere höchstens vom Grade $n-1$, so bestimmen, dass
\weq{3}{\chi(x) = F(x) f(x) + \varphi(x) \psi(x)}
wird, so dass, sobald für $x$ eine Wurzel der Gleichung
\weq{4}{f(x) = 0}
gesetzt wird,
\[
\frac{\chi(x)}{\psi(x)} = \varphi(x). % [sic] period as printed before the following continuation
\]
wird. Da wir nun die Werthe von $y$ nur für solche Werthe von $x$, die Wurzeln von (4) sind, betrachten werden, so verlieren wir nichts an Allgemeinheit, wenn wir von vornherein
\[
y = \varphi(x),
\]
also gleich einer ganzen Function $(n-1)^{\text{ten}}$ Grades setzen. Es sei also
\weq{5}{y = \varphi(x) = \alpha_0 + \alpha_1 x + \alpha_2 x^2 + \cdots + \alpha_{n-1} x^{n-1},}
worin die Co\"efficienten $\alpha$ vorläufig noch ganz unbestimmt bleiben.

Bezeichnen wir mit
\weq{6}{x_1, x_2 \ldots x_n}
%CONT>>

\par\smallskip
{\small 1) Acta eruditorum, Leipzig 1683.}\par
